% Options for packages loaded elsewhere
\PassOptionsToPackage{unicode}{hyperref}
\PassOptionsToPackage{hyphens}{url}
%
\documentclass[
]{article}
\usepackage{amsmath,amssymb}
\usepackage{iftex}
\ifPDFTeX
  \usepackage[T1]{fontenc}
  \usepackage[utf8]{inputenc}
  \usepackage{textcomp} % provide euro and other symbols
\else % if luatex or xetex
  \usepackage{unicode-math} % this also loads fontspec
  \defaultfontfeatures{Scale=MatchLowercase}
  \defaultfontfeatures[\rmfamily]{Ligatures=TeX,Scale=1}
\fi
\usepackage{lmodern}
\ifPDFTeX\else
  % xetex/luatex font selection
\fi
% Use upquote if available, for straight quotes in verbatim environments
\IfFileExists{upquote.sty}{\usepackage{upquote}}{}
\IfFileExists{microtype.sty}{% use microtype if available
  \usepackage[]{microtype}
  \UseMicrotypeSet[protrusion]{basicmath} % disable protrusion for tt fonts
}{}
\makeatletter
\@ifundefined{KOMAClassName}{% if non-KOMA class
  \IfFileExists{parskip.sty}{%
    \usepackage{parskip}
  }{% else
    \setlength{\parindent}{0pt}
    \setlength{\parskip}{6pt plus 2pt minus 1pt}}
}{% if KOMA class
  \KOMAoptions{parskip=half}}
\makeatother
\usepackage{xcolor}
\usepackage{longtable,booktabs,array}
\usepackage{calc} % for calculating minipage widths
% Correct order of tables after \paragraph or \subparagraph
\usepackage{etoolbox}
\makeatletter
\patchcmd\longtable{\par}{\if@noskipsec\mbox{}\fi\par}{}{}
\makeatother
% Allow footnotes in longtable head/foot
\IfFileExists{footnotehyper.sty}{\usepackage{footnotehyper}}{\usepackage{footnote}}
\makesavenoteenv{longtable}
\setlength{\emergencystretch}{3em} % prevent overfull lines
\providecommand{\tightlist}{%
  \setlength{\itemsep}{0pt}\setlength{\parskip}{0pt}}
\setcounter{secnumdepth}{-\maxdimen} % remove section numbering
% unicode-fix.tex — render raw-Unicode symbols that the Latin Modern *text* font
% lacks, using Computer Modern math/text equivalents. This keeps the established
% Computer Modern look while filling the glyphs xelatex would otherwise drop.
% Included via: pandoc ... --include-in-header=unicode-fix.tex
\usepackage{newunicodechar}
\usepackage{pifont}
\usepackage{graphicx} % unicode-math's smallsetminus uses rotatebox
% Use the Computer Modern-compatible set-difference construction; the large
% Unicode glyph selected by TeX Live 2022 is absent from Latin Modern Math.
\AtBeginDocument{\let\setminus\smallsetminus}
% --- Greek letters -> math (Latin Modern Math; consistent with CM body text) ---
\newunicodechar{α}{\ensuremath{\alpha}}
\newunicodechar{β}{\ensuremath{\beta}}
\newunicodechar{γ}{\ensuremath{\gamma}}
\newunicodechar{δ}{\ensuremath{\delta}}
\newunicodechar{ε}{\ensuremath{\varepsilon}}
\newunicodechar{θ}{\ensuremath{\theta}}
\newunicodechar{μ}{\ensuremath{\mu}}
\newunicodechar{ν}{\ensuremath{\nu}}
\newunicodechar{ρ}{\ensuremath{\rho}}
\newunicodechar{σ}{\ensuremath{\sigma}}
\newunicodechar{τ}{\ensuremath{\tau}}
\newunicodechar{φ}{\ensuremath{\varphi}}
\newunicodechar{χ}{\ensuremath{\chi}}
\newunicodechar{ψ}{\ensuremath{\psi}}
\newunicodechar{η}{\ensuremath{\eta}}
\newunicodechar{κ}{\ensuremath{\kappa}}
\newunicodechar{λ}{\ensuremath{\lambda}}
\newunicodechar{π}{\ensuremath{\pi}}
\newunicodechar{ω}{\ensuremath{\omega}}
\newunicodechar{Ḣ}{\ensuremath{\dot H}}
\newunicodechar{≪}{\ensuremath{\ll}}
\newunicodechar{≫}{\ensuremath{\gg}}
\newunicodechar{⟺}{\ensuremath{\Longleftrightarrow}}
\newunicodechar{⟹}{\ensuremath{\Longrightarrow}}
\newunicodechar{⇔}{\ensuremath{\Leftrightarrow}}
\newunicodechar{⟨}{\ensuremath{\langle}}
\newunicodechar{⟩}{\ensuremath{\rangle}}
\newunicodechar{⁰}{\textsuperscript{0}}
\newunicodechar{¹}{\textsuperscript{1}}
\newunicodechar{²}{\textsuperscript{2}}
\newunicodechar{³}{\textsuperscript{3}}
\newunicodechar{⁴}{\textsuperscript{4}}
\newunicodechar{⁵}{\textsuperscript{5}}
\newunicodechar{⁶}{\textsuperscript{6}}
\newunicodechar{⁷}{\textsuperscript{7}}
\newunicodechar{⁸}{\textsuperscript{8}}
\newunicodechar{⁹}{\textsuperscript{9}}
% --- Relations / operators -> math ---
\newunicodechar{↔}{\ensuremath{\leftrightarrow}}
\newunicodechar{⇐}{\ensuremath{\Leftarrow}}
\newunicodechar{⇒}{\ensuremath{\Rightarrow}}
\newunicodechar{∂}{\ensuremath{\partial}}
\newunicodechar{∅}{\ensuremath{\varnothing}}
\newunicodechar{∈}{\ensuremath{\in}}
\newunicodechar{∼}{\ensuremath{\sim}}
\newunicodechar{≈}{\ensuremath{\approx}}
\newunicodechar{≠}{\ensuremath{\neq}}
\newunicodechar{≤}{\ensuremath{\leq}}
\newunicodechar{≥}{\ensuremath{\geq}}
\newunicodechar{⊂}{\ensuremath{\subset}}
\newunicodechar{⊃}{\ensuremath{\supset}}
% --- Special ---
\newunicodechar{ℏ}{\ensuremath{\hbar}}
\newunicodechar{ℤ}{\ensuremath{\mathbb{Z}}}
% --- Superscript modifier letters (footnote-style markers) -> text superscript ---
\newunicodechar{ʰ}{\textsuperscript{h}}
\newunicodechar{ʲ}{\textsuperscript{j}}
\newunicodechar{ᵃ}{\textsuperscript{a}}
\newunicodechar{ᵇ}{\textsuperscript{b}}
\newunicodechar{ᵈ}{\textsuperscript{d}}
\newunicodechar{ᵉ}{\textsuperscript{e}}
\newunicodechar{ᵍ}{\textsuperscript{g}}
\newunicodechar{ᶠ}{\textsuperscript{f}}
\newunicodechar{⁻}{\textsuperscript{\ensuremath{-}}}
% --- Subscript digits -> text subscript ---
\newunicodechar{₀}{\textsubscript{0}}
\newunicodechar{₁}{\textsubscript{1}}
\newunicodechar{₂}{\textsubscript{2}}
\newunicodechar{₃}{\textsubscript{3}}
\newunicodechar{₄}{\textsubscript{4}}
\newunicodechar{₈}{\textsubscript{8}}
% --- Checkmarks ---
\newunicodechar{✓}{\ding{51}}
\newunicodechar{✗}{\ding{55}}

% --- added 2026-08-02 after an independent glyph audit ---
\newunicodechar{ℓ}{\ensuremath{\ell}}
\newunicodechar{⊕}{\ensuremath{\oplus}}
\newunicodechar{⊗}{\ensuremath{\otimes}}
\newunicodechar{∠}{\ensuremath{\angle}}
\newunicodechar{≡}{\ensuremath{\equiv}}
\newunicodechar{∝}{\ensuremath{\propto}}
\newunicodechar{ℝ}{\ensuremath{\mathbb{R}}}
\newunicodechar{₅}{\ensuremath{_5}}
\newunicodechar{₋}{\ensuremath{_-}}
\newunicodechar{⁺}{\ensuremath{^+}}
\newunicodechar{⁽}{\ensuremath{^(}}
\newunicodechar{⁾}{\ensuremath{^)}}
% U+239B..U+23A0: multi-line parenthesis assembly pieces emitted by some
% converters. There is no sane inline equivalent; map to ordinary delimiters so
% the notation survives rather than vanishing silently.
\newunicodechar{⎛}{\ensuremath{\big(}}
\newunicodechar{⎜}{\ensuremath{\big\vert}}
\newunicodechar{⎝}{\ensuremath{\big(}}
\newunicodechar{⎞}{\ensuremath{\big)}}
\newunicodechar{⎟}{\ensuremath{\big\vert}}
\newunicodechar{⎠}{\ensuremath{\big)}}
\ifLuaTeX
  \usepackage{selnolig}  % disable illegal ligatures
\fi
\IfFileExists{bookmark.sty}{\usepackage{bookmark}}{\usepackage{hyperref}}
\IfFileExists{xurl.sty}{\usepackage{xurl}}{} % add URL line breaks if available
\urlstyle{same}
\hypersetup{
  hidelinks,
  pdfcreator={LaTeX via pandoc}}

\author{}
\date{}

\begin{document}

\hypertarget{a-neutrino-mixing-sum-rule-from-octahedral-flavor-structure-2-sinuxb2ux3b8ux2081ux2082-sinuxb2ux3b8ux2082ux2083-76-oux3bbuxb3}{%
\section{A Neutrino Mixing Sum Rule from Octahedral Flavor Structure: 2
sin²θ₁₂ + sin²θ₂₃ = 7/6 +
O(λ³)}\label{a-neutrino-mixing-sum-rule-from-octahedral-flavor-structure-2-sinuxb2ux3b8ux2081ux2082-sinuxb2ux3b8ux2082ux2083-76-oux3bbuxb3}}

\textbf{Author:} Alex Maybaum\\
\textbf{Date:} July 2026

\begin{center}\rule{0.5\linewidth}{0.5pt}\end{center}

\hypertarget{abstract}{%
\subsection{Abstract}\label{abstract}}

Discrete flavor symmetries yield characteristic correlations among the
PMNS mixing angles, with the trimaximal TM1 and TM2 sum rules serving as
benchmarks for reactor experiments. We present a distinct pattern from
octahedral flavor structure: tri-bimaximal mixing at zeroth order,
corrected through O(λ²) at the scale λ = 1/(π√2) = 0.22508 of the cubic
Brillouin-zone construction {[}18{]}. A projection input fixes the
reactor coefficient. A separate normalization b₂₃ = 1 and a
Wilson-coefficient relation (Condition 2), both supplied, give the
leading values sin²θ₁₂ ≈ 0.308003, sin²θ₁₃ ≈ 0.02252, and sin²θ₂₃ ≈
0.55066 without fitting these angles. Condition 2 preserves 2sin²θ₁₂ +
sin²θ₂₃ = 7/6 through O(λ²). The solar value is 0.14σ from the direct
measurement {[}7{]} and 0.07σ from the global fit that includes the
direct data. These overlapping comparisons are retrodictive. The other
angles lie within 1.1σ of the cited NuFIT 6.0 central values. Higher
precision tests the leading values and the angle correlation subject to
open higher-order terms.

\begin{center}\rule{0.5\linewidth}{0.5pt}\end{center}

\hypertarget{introduction}{%
\subsection{1. Introduction}\label{introduction}}

\hypertarget{the-result}{%
\subsubsection{1.1 The result}\label{the-result}}

Tri-bimaximal mixing {[}1{]} was the leading description of the lepton
mixing pattern until Daya Bay established a nonzero θ₁₃ {[}2{]}. The
minimal modified patterns that followed --- TM1 and TM2, each preserving
one column of the TBM matrix --- remain the benchmark correlations
against which reactor-era precision data are analyzed {[}3, 4, 5, 6{]},
and the first JUNO measurement has already sharpened this comparison,
excluding TM2 at more than 3.5σ while leaving TM1 viable {[}9, 10{]}.
Patterns in this class are typically derived from a discrete flavor
symmetry postulated for the lepton sector. In this paper we present a
TBM-deviation pattern of a different origin, with different --- and
sharper --- testable consequences.

The first JUNO measurement of reactor neutrino oscillations {[}7{]}
reports

\textbf{sin²θ₁₂ = 0.3092 ± 0.0087}

from 59.1 days of data. The post-JUNO global fit of Capozzi \emph{et
al.} {[}8{]}, which includes these data and updated SNO+ results, gives

\textbf{sin²θ₁₂ = 0.3085 ± 0.0073}.

Taking the cubic point group O as the lepton flavor structure, with the
normalization and coefficient relation of Sec. 3.4, gives the leading
solar value

\textbf{sin²θ₁₂ = 1/3 − 1/(4π²) ≈ 0.308003} (through O(λ²)).

This lies 0.14σ from the direct measurement {[}7{]} and 0.07σ from the
overlapping global fit without a fit to sin²θ₁₂. The reactor and
atmospheric values are within 1.1σ of the cited NuFIT 6.0 best fits
{[}14{]}. Condition 2 gives the correlation

\textbf{2sin²θ₁₂ + sin²θ₂₃ = 7/6 + O(λ³)}.

The correlation gives a comparison with the TM1, TM2, and TM3
column-preservation patterns. The zeroth-order value 1/3 comes from the
A₄ TBM pattern, and the perturbation uses the Cabibbo scale λ² = 1/(2π²)
of the cubic Brillouin-zone construction.

\hypertarget{where-the-prediction-comes-from}{%
\subsubsection{1.2 Where the prediction comes
from}\label{where-the-prediction-comes-from}}

The perturbation uses the Cabibbo scale λ² = 1/(2π²) (Sec. 2.2), with λ
= 0.22508 compared with 0.22500 ± 0.00067 {[}13{]}. The O(λ²) deviations
are organized by the μ-τ conjugation grading. Condition 2 relates the
U-even and U-odd coefficients of the second-order perturbation, b₂₃ =
(4/3)(b₁₂+b₁₃), with the numerical reading 4/3 = 2A² (Sec. 3.4). The
separate input b₂₃ = 1 fixes the normalization of the solar and
atmospheric shifts. With these inputs and the reactor projection input
(Sec. 3.2), the leading values are

\begin{itemize}
\tightlist
\item
  \textbf{sin²θ₁₂ = 1/3 − 1/(4π²) ≈ 0.308003}
\item
  \textbf{sin²θ₁₃ = 4/(18π²) ≈ 0.02252}
\item
  \textbf{sin²θ₂₃ = 1/2 + 1/(2π²) ≈ 0.55066}
\end{itemize}

The solar comparison constrains the leading pattern with both
coefficient inputs and higher-order contributions open.

\hypertarget{the-empirical-context}{%
\subsubsection{1.3 The empirical context}\label{the-empirical-context}}

The cubic-group flavor structure uses the point group of the simple
cubic lattice. The nearest-neighbor Bravais comparison in {[}18, §4.6
Theorem 7b{]} selects simple cubic within the stated generic-condensate
construction, whose gauge interpretation requires single-copy H-link,
H-cust, and the charged action of {[}18{]}. Its Cabibbo comparison also
uses the momentum-distance prescription of Sec. 2.2. The resulting
flavor comparisons include:

\begin{longtable}[]{@{}
  >{\raggedright\arraybackslash}p{(\columnwidth - 6\tabcolsep) * \real{0.2500}}
  >{\raggedright\arraybackslash}p{(\columnwidth - 6\tabcolsep) * \real{0.2500}}
  >{\raggedright\arraybackslash}p{(\columnwidth - 6\tabcolsep) * \real{0.2500}}
  >{\raggedright\arraybackslash}p{(\columnwidth - 6\tabcolsep) * \real{0.2500}}@{}}
\toprule\noalign{}
\begin{minipage}[b]{\linewidth}\raggedright
Observable
\end{minipage} & \begin{minipage}[b]{\linewidth}\raggedright
Cubic-group prediction
\end{minipage} & \begin{minipage}[b]{\linewidth}\raggedright
Observed
\end{minipage} & \begin{minipage}[b]{\linewidth}\raggedright
Match
\end{minipage} \\
\midrule\noalign{}
\endhead
\bottomrule\noalign{}
\endlastfoot
Cabibbo angle λ & 1/(π√2) = 0.22508 & 0.22500 ± 0.00067 {[}13{]} &
0.04\% \\
Koide angle θ₀ (charged leptons) & C₂/d² = 2/9 & 0.222225 ± 0.000006
{[}13{]} & 0.001\% \\
Solar mixing sin²θ₁₂ (this work, leading) & 1/3 − 1/(4π²) ≈ 0.308003 &
0.3085 ± 0.0073 {[}8{]} & 0.07σ \\
PMNS sum (through O(λ²)) & 2sin²θ₁₂ + sin²θ₂₃ = 7/6 & 1.178 ± 0.020
(Sec. 4.3) & about 0.6σ \\
\end{longtable}

The first two entries are developed in {[}18, §§7.1--7.2{]}. The Cabibbo
geometry and the chirality-based 1/\textbar q\textbar{} power are
structural, with the unit-normalized propagator identification
conditional as in Sec. 2.2. The Koide angle uses the cubic per-axis
normalization, with the amplitude prefactor and operator carrier as
inputs {[}18, §§4.7, 7.2{]}. The solar comparison uses the Layer 2(b)
coefficient relation and a separate normalization (Secs. 3.4 and 6.3).
The sum uses the same solar data and the atmospheric fit specified in
Sec. 4.3.

The first T2K--NOvA joint analysis {[}24{]} has since sharpened
\(|\Delta m^2_{32}|\) below the two-percent level while leaving both the
mass ordering and the \(\theta_{23}\) octant unresolved; the
octant-dependent caveat structure of §7 is unchanged by it.

\hypertarget{position-in-the-broader-literature}{%
\subsubsection{1.4 Position in the broader
literature}\label{position-in-the-broader-literature}}

The lattice construction {[}18{]} belongs to a line of work on the role
of observers in the emergence of quantum theory {[}19, 20, 21, 22{]}.
The flavor calculation in Secs. 2--5 takes the simple-cubic flavor
representation and the stated mixing inputs. Its comparison with data
(Secs. 4--7) uses that pattern independently of the interpretation of
the underlying lattice.

On the Koide side specifically, an independent 2026 construction
{[}23{]} reaches the same two numbers --- the \(\sqrt{2}\) amplitude
from a rank-one coherence condition and \(|\theta_0| = 2/9\) as a
Berry-phase spectral sum \((3/2\pi^2)\sum_{3\nmid n} n^{-2}\) --- by
mechanisms external to the cubic-lattice structural data used here; the
routes' coincidence on both numbers is noted.

The remainder of this paper develops the technical content. Section 2
sets up the cubic-group structure and derives the perturbation scale.
Section 3 computes the deviations from TBM at O(λ²) and derives the sum
rule. Section 4 presents the predictions and compares them with JUNO,
the post-JUNO global fit, and competing TM1/TM2 patterns. Section 5
covers the remaining angles. Section 6 addresses robustness. Section 7
discusses falsifiability under JUNO's projected long-term precision.
Section 8 concludes.

\hypertarget{provenance}{%
\subsubsection{1.5 Provenance}\label{provenance}}

This derivation followed the 2025 JUNO first results {[}7{]} and the
post-JUNO global fit {[}8{]}. The comparison is retrodictive, with no
parameter fitted to either solar value. Future data can test the leading
solar, reactor and atmospheric values and the correlation 2sin²θ₁₂ +
sin²θ₂₃ = 7/6 through O(λ²). A test of the flavor structure also
requires controlling the applicable higher-order terms (Sec. 7).

\hypertarget{cubic-group-flavor-structure}{%
\subsection{2. Cubic-group flavor
structure}\label{cubic-group-flavor-structure}}

\hypertarget{the-cubic-group-o-and-its-aux2084-subgroup}{%
\subsubsection{2.1 The cubic group O and its A₄
subgroup}\label{the-cubic-group-o-and-its-aux2084-subgroup}}

The cubic point group O acts on the six nearest-neighbor signed link
directions \{±ê₁, ±ê₂, ±ê₃\}. Their exact decomposition is 6 = T₁ ⊕ E ⊕
A₁ {[}18, §4.6 Theorem 7{]}. The single-copy H-link premise identifies
these geometric links with one physical carrier. The physical gauge
interpretation of the multiplicities (3, 2, 1) uses H-cust and the
specified matter action {[}18, §§4.5--4.7{]}. The flavor construction
places the three generations in the T₁ sector {[}18, §4.7 Theorem 10{]},
with its physical reading conditional on the carrier and spin premises
there.

The full octahedral rotation group O has order 24 and is isomorphic to
S₄. Its irreducible representations decompose as

\textbf{irreps(O) = A₁ ⊕ A₂ ⊕ E ⊕ T₁ ⊕ T₂}

of dimensions 1, 1, 2, 3, 3. The subgroup A₄ ⊂ O keeps \{A₁, A₂, E, T₁\}
but identifies A₂ with A₁, giving the standard A₄ content A ⊕ A' ⊕ A'\,'
⊕ T of dimensions 1, 1, 1, 3.

We take T₁ as the generation triplet. Realize it by the three Cartesian
unit vectors \{e₁, e₂, e₃\} acted on by the 24 rotations of O. The
democratic vector

\textbf{ĥ = (1/√3)(e₁ + e₂ + e₃)}

is the only direction fixed by the ℤ₃ cyclic subgroup. Its perpendicular
plane is the real two-dimensional complement for that cyclic subgroup
inside T₁. The second TBM column is ĥ, and the first and third lie in
this plane, oriented by μ-τ exchange.

\hypertarget{nearest-neighbor-bravais-comparison}{%
\subsubsection{2.1.1 Nearest-neighbor Bravais
comparison}\label{nearest-neighbor-bravais-comparison}}

Theorem 7b of {[}18, §4.6{]} compares the nearest-neighbor
representations of the fourteen three-dimensional Bravais lattices. In
the generic equivariant-condensate construction of {[}18, §4.5 Theorem
5{]}, the generic unitary stabilizer depends on the condensate's blocks.
The physical gauge interpretation uses single-copy H-link, H-cust, and
the charged action of {[}18{]}.

The non-cubic crystallographic point groups lack the three-dimensional
irreducible block required by this construction. Within the cubic system
the comparison gives:

\begin{longtable}[]{@{}
  >{\raggedright\arraybackslash}p{(\columnwidth - 6\tabcolsep) * \real{0.2500}}
  >{\raggedright\arraybackslash}p{(\columnwidth - 6\tabcolsep) * \real{0.2500}}
  >{\raggedright\arraybackslash}p{(\columnwidth - 6\tabcolsep) * \real{0.2500}}
  >{\raggedright\arraybackslash}p{(\columnwidth - 6\tabcolsep) * \real{0.2500}}@{}}
\toprule\noalign{}
\begin{minipage}[b]{\linewidth}\raggedright
Lattice
\end{minipage} & \begin{minipage}[b]{\linewidth}\raggedright
NN
\end{minipage} & \begin{minipage}[b]{\linewidth}\raggedright
O-decomposition of NN rep
\end{minipage} & \begin{minipage}[b]{\linewidth}\raggedright
Generic unitary condensate stabilizer
\end{minipage} \\
\midrule\noalign{}
\endhead
\bottomrule\noalign{}
\endlastfoot
SC & 6 & T₁ ⊕ E ⊕ A₁ & U(3) × U(2) × U(1) \\
BCC & 8 & A₁ ⊕ A₂ ⊕ T₁ ⊕ T₂ & U(3)² × U(1)² \\
FCC & 12 & A₁ ⊕ E ⊕ T₁ ⊕ 2T₂ & U(3)³ × U(2) × U(1) \\
\end{longtable}

Simple cubic alone supplies the required (3, 2, 1) pattern within this
comparison. Using the momentum-distance prescription of Sec. 2.2 in each
natural triplet sector gives Cabibbo values 0.2251 (SC), 0.1299 (BCC),
and 0.1592 (FCC), compared with 0.22500 ± 0.00067 {[}13{]}. These values
compare the stated nearest-neighbor realizations under that
prescription. The flavor calculation below uses A₄ ⊂ O of the
simple-cubic branch.

\hypertarget{the-cabibbo-scale-from-cubic-brillouin-geometry}{%
\subsubsection{2.2 The Cabibbo scale from cubic Brillouin
geometry}\label{the-cabibbo-scale-from-cubic-brillouin-geometry}}

We identify the perturbation parameter λ with the Cabibbo angle. In a
lattice realization of the cubic-group flavor structure, the three
generations occupy the three T₁ corners of the cubic Brillouin zone at
momenta X\_j = π·e\_j (j = 1, 2, 3). The inter-generation distance in
momentum space is

\textbf{\textbar X\_i − X\_j\textbar{} = π√2 (i ≠ j)}

a fixed geometric constant of the cubic lattice. The leading-order
inter-generation mixing matrix element is the continuum fermion
propagator at this momentum, \textbar M\_ij\textbar{} = \textbar S(X\_j
− X\_i)\textbar{} = 1/\textbar X\_j − X\_i\textbar. The
1/\textbar q\textbar{} scaling (rather than 1/\textbar q\textbar²) is
structural: in the massless limit, the chirality-preserving component of
the fermion propagator at momentum q reduces to S(q) = γ·q/q², and the
taste-changing vertex's Dirac trace contracts this to a magnitude of
order 1/\textbar q\textbar{} rather than the full propagator's
1/\textbar q\textbar². The full derivation, including the chirality
verification of the taste-changing vertex's spinor structure, is given
in {[}18, §7.1{]}. Identifying the matrix element with the Cabibbo angle
gives

\textbf{λ = 1/(π√2) = 0.22508}

matching the observed value λ\_obs = 0.22500 ± 0.00067 {[}13{]} to
0.04\%. Squaring,

\textbf{λ² = 1/(2π²) = 0.05066}

sets the overall scale of O(λ²) corrections. The geometry π√2 and the
chirality-based 1/\textbar q\textbar{} power enter the derivation in
{[}18, §7.1{]}. Identifying the mixing element with the unit-normalized
propagator supplies a separate input, with the emergent dressing open
there.

A second geometric quantity will appear. The angle between any
generation axis e\_i and the democratic direction ĥ satisfies e\_i · ĥ =
1/√3, so cos²(∠(e\_i, ĥ)) = 1/3 and sin²(∠(e\_i, ĥ)) = 2/3. We identify
the Wolfenstein parameter A with this quantity:

\textbf{A = √(2/3) = 0.8165}

agreeing with the PDG value A = 0.826 ± 0.012 at 0.8σ. The combination
A² = 2/3 enters the inter-generation projection geometry, and A⁴ = 4/9
appears directly in the deviation Δ₁₃ below.

\hypertarget{the-pmns-matrix-in-the-altarelli-feruglio-basis}{%
\subsubsection{2.3 The PMNS matrix in the Altarelli-Feruglio
basis}\label{the-pmns-matrix-in-the-altarelli-feruglio-basis}}

We work in the Altarelli-Feruglio (AF) basis {[}12{]}. In this basis the
neutrino sector realizes TBM exactly (after diagonalization of the
seesaw matrix), and all deviations from TBM come from the charged-lepton
mass matrix. The PMNS matrix is

\textbf{U\_PMNS = V\_ℓ† · U\_TBM}

where V\_ℓ is the charged-lepton rotation that diagonalizes M\_ℓ M\_ℓ†,
and

\begin{verbatim}
            ⎛  √(2/3)   √(1/3)    0      ⎞
U_TBM   =   ⎜ -√(1/6)   √(1/3)  -√(1/2)  ⎟
            ⎝ -√(1/6)   √(1/3)   √(1/2)  ⎠
\end{verbatim}

Expanding V\_ℓ in λ,

\textbf{V\_ℓ = exp(λA₁ + λ²A₂ + O(λ³))}

with A₁, A₂ real antisymmetric 3×3 matrices. The remaining task is to
determine A₁ and A₂ from the structural ingredients of the cubic
geometry.

\hypertarget{u-parity-grading-and-the-sum-rule}{%
\subsection{3. U-parity grading and the sum
rule}\label{u-parity-grading-and-the-sum-rule}}

\hypertarget{the-u-parity-generator}{%
\subsubsection{3.1 The U-parity
generator}\label{the-u-parity-generator}}

In the AF basis, the μ-τ exchange U acts on the generation indices by
exchanging the second and third:

\begin{verbatim}
        ⎛ 1  0  0 ⎞
U   =   ⎜ 0  0  1 ⎟
        ⎝ 0  1  0 ⎠
\end{verbatim}

The displayed swap has determinant −1. Its negative −U is a proper cubic
rotation and has the same conjugation action. This action defines a ℤ₂
grading on antisymmetric matrices: A is \textbf{U-even} if UAU⁻¹ = +A
and \textbf{U-odd} if UAU⁻¹ = −A. With (E\_\{ij\})\emph{\{kl\} =
δ}\{ik\}δ\_\{jl\} − δ\_\{il\}δ\_\{jk\}, the decomposition is:

\begin{itemize}
\tightlist
\item
  \textbf{U-even}: E₁₂ + E₁₃
\item
  \textbf{U-odd}: E₁₂ − E₁₃, E₂₃
\end{itemize}

The unique U-odd combination antisymmetric in (μ, τ) is E₂₃; the unique
U-odd combination relating the first generation to the second/third is
E₁₂ − E₁₃.

\hypertarget{determination-of-aux2081}{%
\subsubsection{3.2 Determination of A₁}\label{determination-of-aux2081}}

A₁ is fixed by three structural conditions plus the projection geometry
of §2.2.

The data require sinθ₁₃ = O(λ) (since θ₁₃ ≠ 0), forcing A₁ to contribute
to the (1,3) element of V\_ℓ† U\_TBM. The observation that both Δ₁₂ ≡
sin²θ₁₂ − 1/3 and Δ₂₃ ≡ sin²θ₂₃ − 1/2 are O(λ²) (not O(λ)) further
constrains A₁'s structure. The remaining input is geometric: each end of
an inter-generation mixing vertex projects onto the plane perpendicular
to ĥ, contributing a factor sin(∠(e\_i, ĥ)) = A from the off-democratic
projection (Sec. 2.2). For a single-vertex inter-generation transition,
the total projection factor is A² (one factor per vertex end).
Identifying this projection geometry with the leading O(λ) contribution
to U\_e3 gives the structural identity

\textbf{sinθ₁₃ = A²λ}

at leading O(λ). The coefficient is a supplied projection input that
fixes the overall magnitude of A₁. The O(λ²) amplitude from A₂ enters
the squared reactor angle at higher order (Sec. 3.3).

The result is

\textbf{A₁ = (√2/3)(E₁₂ − E₁₃)}

A₁ is purely U-odd: UA₁U⁻¹ = −A₁. Geometrically, it is a rotation by
angle A² = 2/3 around the axis (e₂ + e₃)/√2, which reproduces sinθ₁₃ =
A²λ as required.

\hypertarget{deviations-to-oux3bbuxb2}{%
\subsubsection{3.3 Deviations to O(λ²)}\label{deviations-to-oux3bbuxb2}}

Expanding V\_ℓ to second order,

\textbf{V\_ℓ = I + λA₁ + λ²(A₂ + ½A₁²) + O(λ³)}

and writing the most general antisymmetric A₂ as

\textbf{A₂ = b₁₂E₁₂ + b₁₃E₁₃ + b₂₃E₂₃}

we compute \textbar U\_PMNS,ij\textbar² to O(λ²). The PMNS angles
satisfy sin²θ₁₃ = \textbar U\_e3\textbar², sin²θ₁₂ cos²θ₁₃ =
\textbar U\_e2\textbar², and sin²θ₂₃ cos²θ₁₃ = \textbar U\_μ3\textbar².
Direct computation (taking V\_ℓ\^{}T = exp(-(λA₁+λ²A₂)) acting on the
columns of U\_TBM, using the explicit form A₁² = diag(-4/9, -2/9, -2/9)
+ (2/9)(e₂⊗e₃ + e₃⊗e₂) for A₁ = (√2/3)(E₁₂−E₁₃)) gives the deviations
from TBM:

\begin{itemize}
\tightlist
\item
  \textbf{Δ₁₃ = (4/9)λ² = A⁴λ²}
\item
  \textbf{Δ₁₂ = −(2/3)(b₁₂ + b₁₃) λ²}
\item
  \textbf{Δ₂₃ = b₂₃ λ²}
\end{itemize}

Three observations follow.

\textbf{(i) A₁ fixes the leading Δ₁₃.} In the TBM sign convention above,
U\_e3 = 2λ/3 + (b₁₂ − b₁₃)λ²/√2 + O(λ³). Its squared modulus through
O(λ²) is (4/9)λ². Interference with the second-order term can contribute
at O(λ³).

\textbf{(ii) The real combination b₁₂ − b₁₃ drops out of the squared
angles through O(λ²).} The mixing matrix is real through the computed
order, so the Dirac Jarlskog invariant vanishes there. Sourcing a
physically complex CP phase remains open.

\textbf{(iii) At O(λ²), Δ₂₃ depends only on b₂₃.} The explicit
calculation gives \textbar U\_μ3\textbar² = ½ + λ²(b₂₃ − 2/9) + O(λ³).
Dividing by cos²θ₁₃ gives sin²θ₂₃ = ½ + λ²b₂₃ + O(λ³).

\hypertarget{sum-rule-and-absolute-normalization}{%
\subsubsection{3.4 Sum rule and absolute
normalization}\label{sum-rule-and-absolute-normalization}}

The deviations depend on two scalar combinations of A₂: (b₁₂ + b₁₃)
controlling Δ₁₂, and b₂₃ controlling Δ₂₃. Their absolute sizes and their
ratio use two inputs.

\textbf{Condition 1 (normalization input).} Take b₂₃ = 1. With λ fixed
independently, this sets the size of the atmospheric correction at O(λ²)
and, with Condition 2, the size of the solar correction.

\textbf{Condition 2 (cubic-group structural relation).} Take

\textbf{b₂₃ = (4/3)(b₁₂ + b₁₃)}

which with Condition 1 gives b₁₂ + b₁₃ = 3/4. Condition 2 fixes the
ratio of the dominant U-odd coefficient to the U-even combination.

\textbf{Leading-order sum rule.} Combining the two deviations from Sec.
3.3 gives

\textbf{2Δ₁₂ + Δ₂₃ = (b₂₃ − (4/3)(b₁₂+b₁₃)) λ² + O(λ³)}.

The coefficient of λ² vanishes precisely when Condition 2 holds. Under
this condition the TBM sum 2sin²θ₁₂ + sin²θ₂₃ = 7/6, with 7/6 = 2/3 +
1/2, is preserved through O(λ²). The expansion leaves contributions at
higher orders open.

\textbf{Structural protection of the sum rule by A₁ alone (Layer 1).}
For the A₁-only rotation V\_ℓ = exp(λA₁), direct calculation gives
2\textbar U\_e2\textbar² + \textbar U\_μ3\textbar² = 7/6 − (14/27)λ² +
O(λ⁴). Dividing by cos²θ₁₃ = 1 − (4/9)λ² + O(λ⁴) cancels the
second-order offset:

**(2sin²θ₁₂ + sin²θ₂₃)\textbar\_\{A₁ only\} = 7/6 + O(λ⁴)**.

The cancellation uses the form of A₁ fixed in Sec. 3.2. The general
expansion also has independent coefficients beyond A₁.

Cond 2 reduces to the question of whether A₂ preserves the sum rule it
was protected for under A₁, or breaks it. The answer depends on the
substrate's specific b coefficients.

\textbf{The factor 4/3.} Cond 2's specific numerical coefficient admits
the structural reading

\textbf{4/3 = 2A² = 2 sin²(∠(e\_i, ĥ))}

where A = √(2/3) is the off-democratic projection of any generation axis
(Sec. 2.2). The identity A² = (d−1)/d = C₂(T₁)/d for d = 3 makes this a
representation-theoretic quantity rather than a numerical coincidence:
it is the squared off-democratic projection for the vector
representation of any rotation group. The factor 4/3 in Condition 2
admits a \emph{partner-count × per-partner projection} reading: the
U-even combination E₁₂+E₁₃ couples generation e₁ to its two partners
\{e₂, e₃\}, with the per-partner projection nominally A². However, the
load-bearing claim of this argument is the \emph{operator alignment} ---
that the substrate's second-order Yukawa structure has a specific T₂
direction matching the partner-count × projection picture. Operator
counting on T₁ ⊗ T₁ → T₂ in the supporting framework {[}18{]} shows that
this T₂ direction is determined by the substrate's specific Yukawa
structure rather than forced by representation theory: the natural
Higgs-democratic spurion does not produce Cond 2.

\textbf{Layered classification of predictions.} Before applying this
terminology to Condition 2, we state operational definitions that
classify predictions by what is required to derive them. The framework's
predictions stratify into four layers:

\begin{itemize}
\tightlist
\item
  \textbf{Layer 0 --- Gauge and discrete structure.} Cubic
  representation theory fixes the six-link decomposition with dimensions
  (3, 2, 1). Single-copy H-link and H-cust with a generic equivariant
  condensate give the unitary stabilizer. The local Standard Model
  algebra additionally uses the specified charged action and gauge
  hypotheses, with the global gauge data requiring the separate inputs
  of {[}18, §4.9{]}. The physical generation and chirality readings use
  H-spin' and H-χ', respectively {[}18, §§4.7--4.8{]}. The flavor
  calculation uses the zeroth-order TBM value sin²θ₁₂ = 1/3.
\item
  \textbf{Layer 1 --- Structural form of corrections.} Properties of the
  \emph{form} of the perturbative expansion that follow from cubic-group
  representation theory and uniqueness constraints. Hold for every
  bijection. Examples: the Cabibbo scale λ² = 1/(2π²) (Brillouin
  geometry, §2.2), the Wolfenstein A² = 2/3 (off-democratic projection),
  sin²θ₁₃ = A⁴λ² (independent of A₂ at O(λ²), as shown in §3.3 above).
\item
  \textbf{Layer 2(a) --- Operator-level structural.} Relations among
  Wilson coefficients of cubic-invariant substrate operators that hold
  for every bijection. Promoted from Layer 2(b) when the coefficient
  ratio reduces to a structural identity for the relevant
  representation.
\item
  \textbf{Layer 2(a)-tractable --- Sharply specified residual.} A
  coefficient ratio that reduces to a sharply specified structural
  quantity (with a representation-theoretic reading) not yet derived
  from substrate dynamics, but with concrete substrate calculations
  available to test it. Distinguished from Layer 2(b) by the sharpness
  of the residual question.
\item
  \textbf{Layer 2(b) --- Solution-specific.} Coefficient values are free
  parameters of the bijection.
\end{itemize}

The \emph{diagnostic} for distinguishing 2(a)-tractable from generic
2(b) is whether the residual question reduces to a \emph{sharply
specified} numerical target --- a single coefficient ratio with a
representation-theoretic reading, testable by definite substrate
calculations --- rather than a generic direction in operator space. A
2(a)-tractable question may still have bijection-specific content; the
label refers to the sharpness of the verification path, not to
representation-theoretic forcing. The identity A² = (d−1)/d = C₂(T₁)/d
for d = 3, by contrast, is forced by representation theory and makes
sin²θ₁₃ = A⁴λ² Layer 1.

\textbf{Layer status of Condition 2.} The A₁ cancellation is a Layer 1
result with the inputs of Sec. 3.2. The ratio b₂₃/(b₁₂+b₁₃) = 4/3 is a
Layer 2(b) input with the geometric reading 4/3 = 2A². The natural
cubic-symmetric Higgs-democratic spurion gives b₁₂:b₁₃:b₂₃ ≈ 17:1:1 in
the heavy-quark hierarchy {[}18, §7.3{]}, leaving a required non-natural
amplitude scaling and channel-specific enhancement (K₂₃/K₁₂ ≈ 2.8
there). The routing calculation excludes the specified second-order
route: its two-hop operator is a spin-taste monomial orthogonal to every
mass channel. The recorded O(λ³) route remains open. The solar and sum
comparisons constrain the proposed leading pattern (Sec. 4).

In standard effective field theory, the U-even and dominant U-odd
coefficients b₁₂+b₁₃ and b₂₃ would be independent Wilson coefficients.
The structural relation b₂₃ = (4/3)(b₁₂+b₁₃) is more naturally
accommodated in frameworks where quantum mechanics is itself an emergent
phenomenon arising from a deeper deterministic structure {[}15, 16,
17{]}, since EFT coefficients can then be related by substrate-level
mechanisms even absent a protecting EFT symmetry. The factorization 4/3
= 2A² is consistent with this picture; whether it follows rigorously
from the cubic-lattice realization of {[}18{]} is the open task
referenced above.

\textbf{Sharpened open task.} In the AF basis (charged leptons diagonal
at LO), second-order perturbative diagonalization of M\_ℓ M\_ℓ†
generates contributions to the A₂ coefficients from cross-terms
involving A₁. Direct computation shows that the framework's choice A₁ =
(√2/3)(E₁₂ − E₁₃) produces a hierarchy-independent cross-term
contribution b₂₃\^{}cross = −1/9 in the m\_e² → 0 limit, while b₁₂ and
b₁₃ receive no analogous cross-term contribution. Cond. 2 at the
observable level then translates to a relation on the substrate's direct
second-order corrections:

\textbf{b₂₃\^{}sub = (4/3)(b₁₂\^{}sub + b₁₃\^{}sub) + 1/9}

The +1/9 term follows from A₁'s perturbative structure under the stated
hierarchy approximation. The required alignment of the direct
corrections is the open task. The cited staggered-fermion calculation
excludes the specified second-order route {[}18, §7.3{]}.

Condition 2 gives 2Δ₁₂ + Δ₂₃ = O(λ³) at the angle level. Section 4.3
compares the leading sum with data. The remaining real coefficient b₁₂ −
b₁₃ has the amplitude and CP scope of Sec. 3.3.

\hypertarget{numerical-predictions}{%
\subsubsection{3.5 Numerical predictions}\label{numerical-predictions}}

Substituting b₂₃ = 1, b₁₂ + b₁₃ = 3/4, and λ² = 1/(2π²) gives the values
through O(λ²):

\begin{itemize}
\tightlist
\item
  \textbf{sin²θ₁₂ = 1/3 − 1/(4π²) ≈ 0.308003}
\item
  \textbf{sin²θ₁₃ = 4/(18π²) ≈ 0.02252}
\item
  \textbf{sin²θ₂₃ = 1/2 + 1/(2π²) ≈ 0.55066}
\end{itemize}

The reactor value uses A₁ and the Cabibbo scale. The solar and
atmospheric values additionally use Conditions 1 and 2. In the minimal
truncation V\_ℓ = exp(λA₁ + λ²A₂), with −0.5 ≤ b₁₂ − b₁₃ ≤ 1 and the
same inputs, keeping higher orders of that exponential shifts sin²θ₁₂ by
+0.0032 to +0.0044. The full expansion has independent O(λ³) terms in
V\_ℓ whose contribution is open.

\hypertarget{the-sinuxb2ux3b8ux2081ux2082-prediction}{%
\subsection{4. The sin²θ₁₂
prediction}\label{the-sinuxb2ux3b8ux2081ux2082-prediction}}

\hypertarget{numerical-value-and-comparison-with-data}{%
\subsubsection{4.1 Numerical value and comparison with
data}\label{numerical-value-and-comparison-with-data}}

The leading solar value from Sec. 3.5 is

\begin{quote}
\textbf{sin²θ₁₂ = 1/3 − 1/(4π²) ≈ 0.308003} (through O(λ²)).
\end{quote}

Comparison with current measurements:

\begin{longtable}[]{@{}lll@{}}
\toprule\noalign{}
Source & sin²θ₁₂ & σ from leading value \\
\midrule\noalign{}
\endhead
\bottomrule\noalign{}
\endlastfoot
This work (leading) & 0.308003 & --- \\
JUNO direct {[}7{]} & 0.3092 ± 0.0087 & 0.14σ \\
Post-JUNO global {[}8{]} & 0.3085 ± 0.0073 & 0.07σ \\
\end{longtable}

JUNO reports the normal-ordering splitting Δm²₂₁ = (7.50 ± 0.12) × 10⁻⁵
eV² {[}7{]}. The global fit including the direct data and updated SNO+
gives Δm²₂₁ = (7.48 ± 0.10) × 10⁻⁵ eV² {[}8{]}. The two solar residuals
therefore use overlapping data.

\hypertarget{comparison-with-column-preservation-patterns}{%
\subsubsection{4.2 Comparison with column-preservation
patterns}\label{comparison-with-column-preservation-patterns}}

Tribimaximal mixing {[}1{]} was the leading phenomenological description
of neutrino oscillation data until Daya Bay measured a nonzero θ₁₃
{[}2{]}. Since then, modified TBM patterns have been the subject of a
substantial literature {[}3, 4, 5, 6{]}. The two best-known
modifications preserve one column of the original TBM matrix exactly:
TM1 (first column) and TM2 (second column). Each gives a one-parameter
family that accommodates the observed θ₁₃ but predicts different values
for the other two angles. Most existing TBM-modification proposals start
from a discrete \emph{flavor symmetry} --- typically A₄, S₄, or a Z₂
residual --- postulated as an axiom of the lepton sector. The minimal
modifications then preserve one column of the TBM matrix or one residual
symmetry of the neutrino mass matrix, yielding one-parameter
correlations between angles. Post-JUNO, He {[}9{]} finds that TM2 is
excluded at \textgreater3.5σ while TM1 remains viable; Zhang {[}10{]}
reports similar results.

Comparison summary:

\begin{longtable}[]{@{}
  >{\raggedright\arraybackslash}p{(\columnwidth - 6\tabcolsep) * \real{0.2632}}
  >{\raggedright\arraybackslash}p{(\columnwidth - 6\tabcolsep) * \real{0.4211}}
  >{\raggedright\arraybackslash}p{(\columnwidth - 6\tabcolsep) * \real{0.2368}}
  >{\raggedright\arraybackslash}p{(\columnwidth - 6\tabcolsep) * \real{0.0789}}@{}}
\toprule\noalign{}
\begin{minipage}[b]{\linewidth}\raggedright
Approach
\end{minipage} & \begin{minipage}[b]{\linewidth}\raggedright
Symmetry origin
\end{minipage} & \begin{minipage}[b]{\linewidth}\raggedright
sin²θ₁₂
\end{minipage} & \begin{minipage}[b]{\linewidth}\raggedright
σ
\end{minipage} \\
\midrule\noalign{}
\endhead
\bottomrule\noalign{}
\endlastfoot
TM1 & postulated A₄ flavor & 0.3184 & 1.35σ \\
TM2 & postulated A₄ flavor & 0.3408 & 4.4σ \\
TM3 & postulated; θ₁₃ = 0 & --- & excluded \\
\textbf{This work (leading)} & \textbf{cubic group O of 3D space} &
\textbf{0.308003} & \textbf{0.07σ} \\
\end{longtable}

(Comparing against post-JUNO global fit 0.3085 ± 0.0073. TM1 uses
sin²θ₁₂ = 1 − (2/3)/cos²θ₁₃; TM2 uses sin²θ₁₂ = 1/(3cos²θ₁₃); both with
sin²θ₁₃ = 0.02195.)

Three differences distinguish the present prediction from TM\_n:

\textbf{(i) Symmetry origin.} The cubic-group calculation uses A₄ ⊂ O on
the T₁ generation triplet of the simple-cubic construction (Sec. 2.1).
The nearest-neighbor gauge and Cabibbo comparisons carry the assumptions
stated there.

\textbf{(ii) Mixing angles.} TM1 gives a correlation between θ₁₂ and
θ₁₃. The cubic-group calculation supplies three leading values under the
inputs of Sec. 3.5 and the leading solar--atmospheric correlation of
Sec. 3.4.

\textbf{(iii) CP violation.} He's surviving pattern predicts sin δ\_CP =
±0.998 {[}9{]}. In the cubic-group calculation the mixing matrix is real
through O(λ²), and a mechanism for a physically complex phase remains
open (Sec. 3.3).

\hypertarget{the-sum-rule-discriminator}{%
\subsubsection{4.3 The sum-rule
discriminator}\label{the-sum-rule-discriminator}}

Condition 2 gives

\textbf{2sin²θ₁₂ + sin²θ₂₃ = 7/6 + O(λ³)}.

The global solar central value 0.3085 and the NuFIT 6.0 normal-ordering
upper-octant value 0.561 give a sum of 1.178. Treating the marginal
errors as uncorrelated and combining them in quadrature gives an
illustrative uncertainty of about 0.020, or a residual of about 0.6σ
from 7/6. The lower-octant central value 0.455 would give about 1.072.
This comparison depends on the atmospheric fit and reuses the solar
data. The leading correlation differs from TM1 and TM2, with the
comparison at higher precision depending on the octant and the
higher-order remainder.

\hypertarget{the-remaining-angles}{%
\subsection{5. The remaining angles}\label{the-remaining-angles}}

\hypertarget{sinuxb2ux3b8ux2081ux2083}{%
\subsubsection{5.1 sin²θ₁₃}\label{sinuxb2ux3b8ux2081ux2083}}

Through O(λ²), the A₂ coefficients drop out of \textbar U\_e3\textbar²
and sin²θ₁₃ = A⁴λ² = 4/(18π²) ≈ 0.02252. The NuFIT 6.0 best fit (normal
ordering) is 0.02195 ± 0.00056, giving a residual of about 1.01σ at the
quoted precision. The leading amplitude sinθ₁₃ = A²λ + O(λ²) gives
sinθ₁₃ ≈ 0.1500 and θ₁₃ ≈ 8.63°, compared with the fitted 8.56° ± 0.10°.

\hypertarget{sinuxb2ux3b8ux2082ux2083}{%
\subsubsection{5.2 sin²θ₂₃}\label{sinuxb2ux3b8ux2082ux2083}}

At leading order, Conditions 1 and 2 give sin²θ₂₃ ≈ 0.5507. The cited
NuFIT 6.0 normal-ordering upper-octant fit is 0.561⁺⁰·⁰¹²₋₀.₀₁₅, giving
a residual of about 0.7σ using the lower marginal error. Condition 2
gives

\textbf{sin²θ₂₃ = 7/6 − 2sin²θ₁₂ + O(λ³)}.

An illustrative ±0.0014 solar uncertainty would propagate to ±0.0028 on
the atmospheric value through this relation. This is the measurement
contribution, with the theory remainder still open.

\hypertarget{summary}{%
\subsubsection{5.3 Summary}\label{summary}}

\begin{longtable}[]{@{}llll@{}}
\toprule\noalign{}
Angle & Leading value & Measurement & σ \\
\midrule\noalign{}
\endhead
\bottomrule\noalign{}
\endlastfoot
sin²θ₁₂ & 1/3 − 1/(4π²) ≈ 0.308003 & 0.3085 ± 0.0073 & 0.07 \\
sin²θ₁₃ & 4/(18π²) = 0.02252 & 0.02195 ± 0.00056 & 1.01 \\
sin²θ₂₃ & 1/2 + 1/(2π²) = 0.5507 & 0.561⁺⁰·⁰¹²₋₀.₀₁₅ & 0.7 \\
\end{longtable}

The three leading values lie within 1.1σ of these fitted central values.

\hypertarget{robustness}{%
\subsection{6. Robustness}\label{robustness}}

\hypertarget{cumulative-empirical-match}{%
\subsubsection{6.1 Cumulative empirical
match}\label{cumulative-empirical-match}}

The flavor comparisons use cubic-group inputs in several sectors. The
sum comparison uses the angle data in Sec. 4.3.

\begin{longtable}[]{@{}
  >{\raggedright\arraybackslash}p{(\columnwidth - 8\tabcolsep) * \real{0.2000}}
  >{\raggedright\arraybackslash}p{(\columnwidth - 8\tabcolsep) * \real{0.2000}}
  >{\raggedright\arraybackslash}p{(\columnwidth - 8\tabcolsep) * \real{0.2000}}
  >{\raggedright\arraybackslash}p{(\columnwidth - 8\tabcolsep) * \real{0.2000}}
  >{\raggedright\arraybackslash}p{(\columnwidth - 8\tabcolsep) * \real{0.2000}}@{}}
\toprule\noalign{}
\begin{minipage}[b]{\linewidth}\raggedright
Observable
\end{minipage} & \begin{minipage}[b]{\linewidth}\raggedright
Structural input
\end{minipage} & \begin{minipage}[b]{\linewidth}\raggedright
Prediction
\end{minipage} & \begin{minipage}[b]{\linewidth}\raggedright
Observed
\end{minipage} & \begin{minipage}[b]{\linewidth}\raggedright
Match
\end{minipage} \\
\midrule\noalign{}
\endhead
\bottomrule\noalign{}
\endlastfoot
Cabibbo angle & Brillouin distance π√2; chirality 1/\textbar q\textbar{}
& λ = 1/(π√2) = 0.22508 & 0.22500 ± 0.00067 & 0.04\% \\
Koide angle (charged leptons) & C₂/d² for cubic group & θ₀ = 2/9 =
0.22222 & 0.222225 ± 0.000006 & 0.001\% \\
Solar mixing sin²θ₁₂ (leading) & TBM, Cabibbo scale, Conditions 1 and 2
& 1/3 − 1/(4π²) ≈ 0.308003 & 0.3085 ± 0.0073 & 0.07σ \\
PMNS sum (through O(λ²)) & Condition 2 & 2sin²θ₁₂ + sin²θ₂₃ = 7/6 &
1.178 ± 0.020 (Sec. 4.3) & about 0.6σ \\
Reactor mixing sin²θ₁₃ (leading) & A²λ projection input & 4/(18π²) ≈
0.02252 & 0.02195 ± 0.00056 & 1.01σ \\
Atmospheric mixing sin²θ₂₃ (leading) & TBM, Cabibbo scale, Condition 1 &
1/2 + 1/(2π²) ≈ 0.55066 & 0.561⁺⁰·⁰¹²₋₀·₀₁₅ & about 0.7σ \\
\end{longtable}

The Cabibbo scale λ² and the projection factor A² enter several rows.
The solar residuals use overlapping datasets, and the sum row is a
function of the solar and atmospheric fits.

\hypertarget{sensitivity-to-condition-2}{%
\subsubsection{6.2 Sensitivity to Condition
2}\label{sensitivity-to-condition-2}}

Condition 2 uses the numerical identity 4/3 = 2A², with A² = 2/3 from
projection perpendicular to ĥ. The identification of 4/3 with the
Wilson-coefficient ratio is a separate input.

Holding b₂₃ and λ fixed and neglecting higher orders, a fractional
coefficient shift 4/3 → (4/3)(1 + ε) changes sin²θ₁₂ by εΔ₂₃/2 to first
order. With Δ₂₃ = 1/(2π²), a 10\% shift gives about 0.0025, about 29\%
of the direct first-result uncertainty. Under these conditions the
direct, global, and illustrative future solar uncertainties of 0.0087,
0.0073, and 0.0014 correspond to about 34\%, 29\%, and 5.5\%,
respectively, on the coefficient. Allowing higher orders or varying the
normalization changes that sensitivity.

\hypertarget{status-of-condition-2}{%
\subsubsection{6.3 Status of Condition 2}\label{status-of-condition-2}}

Condition 2 has a Layer 1 component, the A₁ cancellation of Sec. 3.4,
and a Layer 2(b) input specifying b₂₃/(b₁₂+b₁₃) = 4/3. The reactor angle
at this order depends on A₁ and inherits the open emergent dressing of
the Cabibbo identification {[}18, §7.1{]}. The absolute solar and
atmospheric shifts use the separate normalization b₂₃ = 1.

The staggered-fermion calculation {[}18, §7.3{]} excludes the specified
two-hop route to the ratio because the resulting spin-taste monomial is
orthogonal to every mass channel. The candidate bijections in {[}18{]}
use linear dynamics (A5), or the separate rounded realization studied
there, with T-invariance and P-indivisibility at their stated scopes.
Whether the rounded bijection samples any energy shell is open. The
tested quadratic-energy level sets are not invariant under that map. The
exact linear rule has energy shells larger than its order that are not
single cycles at prime q not dividing the lattice size ({[}18, §4.2{]}).

The natural cubic-symmetric staggered Yukawa structure gives b₁₂:b₁₃:b₂₃
≈ 17:1:1 in the stated heavy-quark hierarchy {[}18, §7.3{]}. This result
and the two-hop exclusion constrain the candidate constructions. The
solar and sum residuals compare the proposed leading pattern with
observations. A bijection supplying the required alignment remains open.

\hypertarget{falsifiability}{%
\subsection{7. Falsifiability}\label{falsifiability}}

\textbf{Solar value.} At an illustrative solar uncertainty of ±0.0014,
unchanged direct and global central values would lie about 0.855σ and
0.355σ from 0.308003, respectively. This tests a model taking the O(λ²)
value as its complete solar prediction. A comparison testing the
proposed flavor structure also requires the normalization and
coefficient inputs and a controlled higher-order remainder. The shifts
in Sec. 3.5 apply to the specified minimal truncation.

\textbf{Solar--atmospheric sum.} The 7/6 value holds through O(λ²) under
Condition 2. Solar and atmospheric measurements can test a model
imposing the exact sum as an additional phenomenological commitment. A
test of Condition 2 itself uses the octant and the experimental
covariance as applicable and requires control of the O(λ³) remainder.

\textbf{Reactor angle and CP.} Reactor measurements test the leading
value from A₁ and λ subject to higher orders. Long-baseline experiments
can constrain δ\_CP. The present real calculation supplies a vanishing
Dirac invariant at its computed order and leaves the source of a complex
phase open.

\hypertarget{conclusion}{%
\subsection{8. Conclusion}\label{conclusion}}

The octahedral flavor calculation gives the leading solar value sin²θ₁₂
= 1/3 − 1/(4π²) ≈ 0.308003 under the separate normalization b₂₃ = 1 and
Condition 2. It lies 0.14σ from the direct measurement {[}7{]} and 0.07σ
from the overlapping global fit, without fitting either measurement. The
leading reactor and atmospheric values are within 1.1σ of the cited
NuFIT 6.0 best fits. The reactor calculation uses the A₁ projection
input and inherits the open emergent dressing of λ. The solar and
atmospheric correlation 2sin²θ₁₂ + sin²θ₂₃ = 7/6 + O(λ³) follows from
Condition 2, whose source is a Layer 2(b) question with the specified
second-order route excluded {[}18, §7.3{]}.

Future reactor and atmospheric precision can test the leading values and
correlation. A test of the proposed flavor structure uses the stated
coefficient inputs, octant, and higher-order remainder.

\hypertarget{acknowledgments}{%
\subsection{Acknowledgments}\label{acknowledgments}}

The author thanks the JUNO Collaboration for the rapid public release of
their first oscillation results, and F. Capozzi, E. Lisi, F. Marcone, A.
Marrone, and A. Palazzo for the prompt post-JUNO global fit that enabled
the present comparison. Discussions with the broader
emergent-quantum-mechanics community have informed the framework's
positioning relative to existing approaches; specific intellectual debts
are recorded in the references and in the underlying framework
documentation {[}18{]}.

AI assistance in preparing this work is disclosed at\\
\url{https://github.com/amaybaum/incompleteness\#ai-assistance}.

\hypertarget{references}{%
\subsection{References}\label{references}}

{[}1{]} P. F. Harrison, D. H. Perkins, W. G. Scott, ``Tri-bimaximal
mixing and the neutrino oscillation data,'' Phys. Lett. B 530, 167
(2002).

{[}2{]} F. P. An et al.~(Daya Bay), ``Observation of
electron-antineutrino disappearance at Daya Bay,'' Phys. Rev.~Lett. 108,
171803 (2012).

{[}3{]} S. F. King and C. Luhn, ``Neutrino mass and mixing with discrete
symmetry,'' Rep.~Prog. Phys. 76, 056201 (2013).

{[}4{]} X.-G. He and A. Zee, ``Minimal modification to tribimaximal
mixing,'' Phys. Rev.~D 84, 053004 (2011).

{[}5{]} X.-G. He, ``TM1 mixing pattern from A₄,'' Chin. J. Phys. 53,
100101 (2015).

{[}6{]} S. T. Petcov and A. V. Titov, ``Assessing the viability of A₄,
S₄ and A₅ flavour symmetries for description of neutrino mixing,'' Phys.
Rev.~D 97, 115045 (2018).

{[}7{]} JUNO Collaboration (A. Abusleme et al.), ``First measurement of
reactor neutrino oscillations at JUNO,'' arXiv:2511.14593 (2025); Nature
654, 343 (2026).

{[}8{]} F. Capozzi, E. Lisi, F. Marcone, A. Marrone, A. Palazzo,
``Updated bounds on the (1,2) neutrino oscillation parameters after
first JUNO results,'' arXiv:2511.21650 (2025); Phys. Rev.~D 114, 016026
(2026).

{[}9{]} X.-G. He, ``Modified Tri-bimaximal neutrino mixing confronted by
JUNO θ₁₂ measurement,'' arXiv:2511.15978 (2025).

{[}10{]} D. Zhang, ``Trimaximal mixing patterns meet the first JUNO
result,'' arXiv:2511.15654 (2025).

{[}11{]} E. Ma and G. Rajasekaran, ``Softly broken A₄ symmetry for
nearly degenerate neutrino masses,'' Phys. Rev.~D 64, 113012 (2001).

{[}12{]} G. Altarelli and F. Feruglio, ``Discrete flavor symmetries and
models of neutrino mixing,'' Rev.~Mod. Phys. 82, 2701 (2010).

{[}13{]} S. Navas et al.~(Particle Data Group), Phys. Rev.~D 110, 030001
(2024).

{[}14{]} I. Esteban, M. C. Gonzalez-Garcia, M. Maltoni, I.
Martinez-Soler, J. P. Pinheiro, T. Schwetz, ``NuFit-6.0: updated global
analysis of three-flavor neutrino oscillations,'' JHEP 12 (2024) 216;
www.nu-fit.org.

{[}15{]} G. 't Hooft, \emph{The Cellular Automaton Interpretation of
Quantum Mechanics}, Fundamental Theories of Physics, vol.~185 (Springer,
2016); arXiv:1405.1548.

{[}16{]} S. L. Adler, \emph{Quantum Theory as an Emergent Phenomenon:
The Statistical Mechanics of Matrix Models as the Precursor of Quantum
Field Theory} (Cambridge University Press, 2004).

{[}17{]} J. A. Barandes, ``The Stochastic-Quantum Correspondence,''
arXiv:2302.10778 (2023).

{[}18{]} A. Maybaum, \emph{The Incompleteness of Observation} (Zenodo
archive, 2026), https://doi.org/10.5281/zenodo.19060318.

{[}19{]} V. Chandrasekaran, R. Longo, G. Penington, E. Witten, ``An
algebra of observables for de Sitter space,'' \emph{JHEP} \textbf{02}
(2023) 082; arXiv:2206.10780.

{[}20{]} J. Maldacena, ``Real observers solving imaginary problems,''
arXiv:2412.14014 (2024).

{[}21{]} D. Harlow, M. Usatyuk, Y. Zhao, ``Quantum mechanics and
observers for gravity in a closed universe,'' \emph{JHEP} \textbf{02}
(2026) 108; arXiv:2501.02359.

{[}22{]} K. Slagle and J. Preskill, ``Emergent Quantum Mechanics at the
Boundary of a Local Classical Lattice Model,'' \emph{Phys. Rev.~A}
\textbf{108}, 012217 (2023); arXiv:2207.09465.

{[}23{]} K. Shulga, ``Charged-Lepton Koide Geometry from a Green-Dressed
Compact Family Cycle,'' arXiv:2605.10245 (2026).

{[}24{]} NOvA and T2K Collaborations (S. Abubakar et al.), ``Joint
neutrino oscillation analysis from the T2K and NOvA experiments,''
Nature 646, 818 (2025); arXiv:2510.19888.

\end{document}
% source-sha256: a589243f04856fed9db9bd5199be4f3c50e68b49383475596576e238c68fc974
